\section{SpeechGR}
\label{sec:method}

\paragraph{Task formulation.}
Let \(\mathcal{C}=\{(p_i,y_i)\}_{i=1}^{N}\) be a corpus of spoken
passages, where \(p_i\) is an audio segment and \(y_i\) is its unique
document identifier (DocID). Let
\(\mathcal{Q}=\{(q_j,t(j))\}_{j=1}^{M}\) map each spoken query \(q_j\) to its
gold passage \(p_{t(j)}\). SpeechGR learns an encoder--decoder model
\(P_\theta(y\mid \mathbf{u}(q))\) and ranks passages by the probability of their
identifiers. Both queries and passages are represented as acoustic tokens.
The SLUE-SQA-5 relevance relation differs from IRGen's same-modality image
task: the gold passage is labelled because its transcript contains information
that answers the question, not because the query and passage recordings are
acoustically similar. SpeechGR is therefore trained to predict a directed
question-to-passage association rather than an acoustically similar recording.
Because SLUE-SQA-5 is text-derived, the experiment cannot determine how much
of this association reflects semantic reasoning rather than lexical
correspondence. Unlike IRGen's learned visual codes, SpeechGR generates
corpus-provided metadata DocIDs; their question-to-passage associations are
learned through the joint indexing and retrieval paths below.

We use \emph{end-to-end} for the ASR-free retrieval path from speech units to
DocIDs, not for joint optimization of the frozen HuBERT--\(k\)-means frontend
with the retriever. We use \emph{textless} for task inputs: SpeechGR's
acoustic-unit adaptation, retrieval fine-tuning, and inference do not consume
query or passage transcripts or a natural-language retrieval instruction.
Flan-T5 is nevertheless text-pretrained, and the DocIDs are metadata-derived
strings. We call this setup
\emph{index-free} in the standard generative-retrieval sense: inference does
not store passage embeddings or perform approximate nearest-neighbour (ANN)
search.  The valid-DocID trie used during decoding is an output constraint,
not an embedding index or similarity-search structure.

\begin{figure}[t]
    \centering
    \resizebox{\linewidth}{!}{%
    \begin{tikzpicture}[
        font=\small,
        >=Latex,
        flow/.style={-{Latex[length=2.2mm]}, line width=0.7pt},
        audio/.style={draw=blue!65!black, fill=blue!8, rounded corners=2pt,
            minimum height=7mm, align=center},
        units/.style={draw=teal!65!black, fill=teal!8, rounded corners=2pt,
            minimum height=7mm, align=center},
        model/.style={draw=orange!75!black, fill=orange!10, rounded corners=2pt,
            minimum height=9mm, align=center},
        output/.style={draw=red!65!black, fill=red!7, rounded corners=2pt,
            minimum height=7mm, align=center},
        constraint/.style={draw=purple!65!black, fill=purple!7,
            rounded corners=2pt, minimum height=7mm, align=center}
    ]
        % Shared acoustic tokenizer.
        \node[audio, minimum width=18mm] (wave) at (0,3.4) {waveform\\\(a\)};
        \node[audio, minimum width=22mm] (hubert) at (2.5,3.4) {HuBERT\\layer \(\ell\)};
        \node[units, minimum width=22mm] (kmeans) at (5.2,3.4) {\(k\)-means\\codebook \(K\)};
        \node[units, minimum width=22mm] (dedup) at (7.9,3.4) {collapse\\repetitions};
        \node[units, minimum width=18mm] (useq) at (10.4,3.4) {units\\\(\mathbf{u}(a)\)};
        \draw[flow] (wave) -- (hubert);
        \draw[flow] (hubert) -- (kmeans);
        \draw[flow] (kmeans) -- (dedup);
        \draw[flow] (dedup) -- (useq);

        % Fine-tuning: indexing and retrieval examples share parameters.
        \node[audio, minimum width=17mm] (passage) at (0,1.75)
            {passage \(p_i\)};
        \node[units, minimum width=16mm] (pu) at (2.0,1.75)
            {\(\mathbf{u}(p_i)\)};
        \node[audio, minimum width=17mm] (query) at (0,0.65)
            {question \(q_j\)};
        \node[units, minimum width=16mm] (qu) at (2.0,0.65)
            {\(\mathbf{u}(q_j)\)};
        \node[model, minimum width=21mm] (t5) at (4.3,1.2)
            {shared\\Flan-T5};
        \node[output, minimum width=20mm] (target) at (6.7,1.2)
            {target DocID\\\(y_i\) or \(y_{t(j)}\)};
        \draw[flow] (passage) -- (pu);
        \draw[flow] (query) -- (qu);
        \draw[flow] (pu) -- (t5);
        \draw[flow] (qu) -- (t5);
        \draw[flow] (t5) -- (target);

        % Inference: the trie constrains decoder scores, not a gold target.
        \node[units, minimum width=16mm] (infq) at (8.8,1.2)
            {\(\mathbf{u}(q)\)};
        \node[model, minimum width=21mm] (decoder) at (10.9,1.2)
            {Flan-T5\\decoder};
        \node[output, minimum width=23mm] (generated) at (13.5,1.2)
            {generated\\DocID};
        \node[constraint, minimum width=22mm] (trie) at (10.9,0.05)
            {valid-DocID trie};
        \draw[flow] (infq) -- (decoder);
        \draw[flow] (decoder) -- (generated);
        \draw[flow, dashed] (trie) -- node[right, font=\scriptsize]
            {valid next tokens} (decoder);

        \begin{scope}[on background layer]
            \node[draw=gray!55, rounded corners=3pt, inner sep=5pt,
                fit=(wave)(hubert)(kmeans)(dedup)(useq),
                label={[font=\scriptsize]above:shared acoustic tokenizer}] {};
            \node[draw=gray!55, rounded corners=3pt, inner sep=5pt,
                fit=(passage)(query)(pu)(qu)(t5)(target),
                label={[font=\scriptsize]above:fine-tuning}] {};
            \node[draw=gray!55, rounded corners=3pt, inner sep=5pt,
                fit=(infq)(decoder)(generated)(trie),
                label={[font=\scriptsize]above:inference}] {};
        \end{scope}
    \end{tikzpicture}%
    }
    \caption{\textbf{SpeechGR overview.} A shared HuBERT--\(k\)-means
    preprocessor maps passages and queries to deduplicated acoustic units.
    Fine-tuning combines passage-to-DocID indexing examples and
    spoken-question-to-answer-passage-DocID retrieval examples. At inference, a trie masks
    invalid next-token continuations during beam search; it stores valid output
    prefixes, not document embeddings.}
    \label{fig:speechgr}
\end{figure}


\subsection{Discrete speech units}
\label{sec:units}

Given an audio waveform \(a\), a pretrained HuBERT encoder
\citep{Hsu2021HuBERTSS} produces frame representations
\(\mathbf{z}_{1:T}\), with \(\mathbf{z}_t\in\mathbb{R}^{d}\).  We quantize
each frame using a \(k\)-means codebook
\(\mathcal{K}=\{\mathbf{c}_1,\ldots,\mathbf{c}_K\}\):
\begin{equation}
    u_t=\arg\min_{j\in\{1,\ldots,K\}}
    \lVert\mathbf{z}_t-\mathbf{c}_j\rVert_2^2 .
    \label{eq:quantization}
\end{equation}
This yields the discrete sequence
\(\mathbf{u}(a)=(u_1,\ldots,u_T)\).  Consecutive repetitions are collapsed,
so a run \(u_t=\cdots=u_{t+r}\) contributes one token.  This removes much of
the frame-rate redundancy without requiring a transcript.  We also test byte
pair encoding (BPE), following prior acoustic-unit work
\citep{Wang2024ACS}, as an optional second compression stage that merges
frequent adjacent units.  BPE is an ablation rather than part of the final
SpeechGR pipeline. Section~\ref{sec:analysis} reports the two same-\(K\)
comparisons available in the experiment record. The final reported
configuration uses HuBERT-large layer~22, \(K=500\), and deduplication without
BPE.

\subsection{Identifiers and model adaptation}
\label{sec:identifier-model}

The main identifier design uses the corpus-provided DocIDs.  These are short
character strings that encode corpus metadata, such as an article prefix and
sentence offset, and are tokenized into output sequences by the model
tokenizer.  We additionally study unstructured atomic identifiers, assigning
one distinct identifier to each passage, as an ablation motivated by
non-text generative retrieval \citep{li-etal-2024-generative}.  The main
system retains the corpus-provided string identifiers because they give the
better Hit@20 in our experiments.

SpeechGR uses Flan-T5-base \citep{10.5555/3722577.3722647} as its
encoder--decoder backbone.  To accept the acoustic sequence
\(\mathbf{u}(a)\), we expand its vocabulary with one entry per discrete unit.
Each added embedding is initialized by copying the vector of an unused token
in the original text vocabulary and remains trainable thereafter.  The
pretrained text vocabulary is retained for sentinel symbols and DocID
generation.

\subsection{Speech-unit pretraining}
\label{sec:speech-pretraining}

Before retrieval training, we adapt the expanded model to acoustic-token
sequences using the T5 span-corruption objective
\citep{Raffel2019ExploringTL}.  On a 6,000-hour subset of LibriLight, 15\% of
the units are grouped into contiguous spans and replaced by distinct sentinel
tokens; the decoder reconstructs the removed spans in order.  The final model
uses ten epochs of this pretraining.  Importantly, this stage adapts the
expanded encoder--decoder to discrete speech units; it does
\emph{not} teach the model the target Spoken Wikipedia passage-to-DocID
mapping.  Those corpus associations are learned only in the next stage.

\subsection{Joint indexing and retrieval training}
\label{sec:training}

Retrieval training combines two kinds of sequence-to-sequence examples,
following the two roles of a generative retriever
\citep{Tay2022TransformerMA}.  The \emph{indexing path} maps each spoken
passage to its own identifier,
\[
    \mathbf{u}(p_i)\longrightarrow y_i ,
\]
thereby associating corpus content with a DocID in the model parameters.  The
\emph{retrieval path} maps a spoken query to the identifier of its labelled
passage,
\[
    \mathbf{u}(q_j)\longrightarrow y_{t(j)} .
\]
Both paths use the same discrete-unit preprocessing and update the shared
Flan-T5 encoder--decoder on longer passage inputs and shorter query inputs.
For either training pair
\((\mathbf{x},\mathbf{y})\), where
\(\mathbf{y}=(y_1,\ldots,y_L)\), SpeechGR minimizes token-level
cross-entropy:
\begin{equation}
    \mathcal{L}_{\mathrm{CE}}(\mathbf{x},\mathbf{y})
    =-\sum_{\ell=1}^{L}
    \log P_\theta(y_\ell\mid y_{<\ell},\mathbf{x}) .
    \label{eq:ce}
\end{equation}
The paper's ranking-loss and windowed Q-Former variants are evaluated
separately in Section~\ref{sec:analysis}; neither is used in the final method.

\subsection{Constrained retrieval}
\label{sec:decoding}

At inference, SpeechGR converts only the spoken query to discrete units and
generates a DocID autoregressively.  We build a trie from the tokenized DocIDs
of all passages in \(\mathcal{C}\).  At every decoding step, tokens that do not
extend the current prefix along a trie edge are masked; completed hypotheses
therefore correspond to valid corpus passages.  Beam search retains the 20
highest-scoring prefixes, and completed beams give the ranked retrieval list.
This constraint prevents invalid identifiers but does not compare the query
against stored passage vectors, preserving the index-free GR formulation
defined above.
